\uuid{OMpD}
\titre{Structures non usuelles}
\auteur{hanani ; mzari}
\organisation{exo7}
\licence{CC BY-NC-SA 4.0}
\datecreate{2020-01-20}
\niveau{L1}
\module{Algèbre}
\chapitre{Espace vectoriel}
\sousChapitre{Définition, sous-espace}
\difficulte{3}
\theme{Espaces vectoriels, Structures non usuelles}
\exo7qcmid{lille-2.espaces-vectoriels.022}

\contenu{
\begin{question}
\qclassification{106.01}
\qtags{motcle=Structures non usuelles}
On considère $E= (\Rr^*)^2$ muni de l'addition et la multiplication par un réel suivantes :
$$(x,y)+(x',y')=(xx',yy')\quad \mbox{et}\quad \lambda .(x,y)= (\lambda x,\lambda y).$$
Quelles sont les assertions vraies ?
\begin{answers}
\bad{$E$ est stable par multiplication par un scalaire.}
\bad{L'élément neutre pour l'addition est $(0,0)$.}
\good{ L'inverse, pour l'addition, de $(x,y)$ est $\displaystyle \left(\frac{1}{x}, \frac{1}{y}\right)$.}
\bad{$E$ est un $\Rr$-espace vectoriel.}
\end{answers}
\begin{explanations} On vérifie que l'élément neutre pour l'addition est $(1,1)$, que l'inverse pour l'addition d'un couple $(x,y) \in (\Rr^*)^2$ est $\displaystyle\left(\frac{1}{x}, \frac{1}{y}\right)$ et que $E$ n'est pas stable par multiplication par le scalaire zéro. Par conséquent, $E$ n'est pas un espace vectoriel.
\end{explanations}
\end{question}
}
